\documentclass[aspectratio=169,11pt]{beamer}

\usetheme{Madrid}
\usecolortheme{default}

\usepackage[T1]{fontenc}
\usepackage[utf8]{inputenc}
\usepackage[lithuanian]{babel}
\usepackage{amsmath,amssymb}

\title{Pilnojo kvadrato išskyrimas}
\author{Andrius Berniukevičius}
\date{}

\begin{document}

\begin{frame}
    \titlepage
\end{frame}

\begin{frame}{Kas yra pilnasis kvadratas?}

\begin{block}{Apibrėžimas}
Pilnasis kvadratas -- tai reiškinys, kurį galima užrašyti dvinario kvadratu:
\[
(a+b)^2 \quad \text{arba} \quad (a-b)^2.
\]
\end{block}

\pause

\begin{block}{Pilnojo kvadrato išskyrimas}
Pilnojo kvadrato išskyrimas -- tai trinario pertvarkymas į dvinario
kvadratą arba į dvinario kvadrato ir skaičiaus sumą ar skirtumą.
\end{block}

\end{frame}

\begin{frame}{Pilnojo kvadrato trinaris}

Jei trinaris yra formos
\[
a^2+2ab+b^2,
\]
\pause
tai
\[
a^2+2ab+b^2=(a+b)^2.
\]

\pause

Jei trinaris yra formos
\[
a^2-2ab+b^2,
\]
\pause
tai
\[
a^2-2ab+b^2=(a-b)^2.
\]

\end{frame}

\begin{frame}{Kaip atpažinti pilnojo kvadrato trinarį?}

\begin{enumerate}[<+->]
    \item Patikriname, ar pirmasis narys yra kvadratas.
    \item Patikriname, ar paskutinis narys yra kvadratas.
    \item Patikriname, ar vidurinis narys lygus
    \[
    2ab \quad \text{arba} \quad -2ab.
    \]
\end{enumerate}

\end{frame}

\begin{frame}{Pavyzdys 1}

Išskirkite pilnąjį kvadratą:
\[
x^2+6x+9.
\]

\pause

Pirmasis narys yra $x^2$, todėl $a=x$.

\pause

Paskutinis narys yra $9=3^2$, todėl $b=3$.

\pause

Patikriname vidurinį narį:
\[
2ab=2\cdot x\cdot3=6x.
\]

\pause

Todėl
\[
x^2+6x+9=x^2+2\cdot x\cdot3+3^2=(x+3)^2.
\]

\end{frame}

\begin{frame}{Pavyzdys 2}

Išskirkite pilnąjį kvadratą:
\[
4x^2-12x+9.
\]

\pause

Pirmasis ir paskutinis nariai yra kvadratai:
\[
4x^2=(2x)^2,
\qquad
9=3^2.
\]

\pause

Patikriname vidurinį narį:
\[
-2\cdot2x\cdot3=-12x.
\]

\pause

Todėl
\[
4x^2-12x+9
=(2x)^2-2\cdot2x\cdot3+3^2
=(2x-3)^2.
\]

\end{frame}

\begin{frame}{Pavyzdys 3}

Išskirkite pilnąjį kvadratą:
\[
9a^2+24ab+16b^2.
\]

\pause

Pirmasis ir paskutinis nariai yra kvadratai:
\[
9a^2=(3a)^2,
\qquad
16b^2=(4b)^2.
\]

\pause

Patikriname vidurinį narį:
\[
2\cdot3a\cdot4b=24ab.
\]

\pause

Todėl
\[
9a^2+24ab+16b^2
=(3a)^2+2\cdot3a\cdot4b+(4b)^2
=(3a+4b)^2.
\]

\end{frame}

\begin{frame}{Pavyzdys 4}

Išskirkite pilnąjį kvadratą:
\[
4a^2+12ab+9b^2.
\]

\pause

Pirmasis ir paskutinis nariai yra kvadratai:
\[
4a^2=(2a)^2,
\qquad
9b^2=(3b)^2.
\]

\pause

Patikriname vidurinį narį:
\[
2\cdot2a\cdot3b=12ab.
\]

\pause

Todėl
\[
4a^2+12ab+9b^2
=(2a)^2+2\cdot2a\cdot3b+(3b)^2
=(2a+3b)^2.
\]

\end{frame}

\begin{frame}{Pavyzdys 5}

Išskirkite pilnąjį kvadratą:
\[
9m^2n^2-6mny+y^2.
\]

\pause

Pirmasis ir paskutinis nariai yra kvadratai:
\[
9m^2n^2=(3mn)^2,
\qquad
y^2=y^2.
\]

\pause

Patikriname vidurinį narį:
\[
-2\cdot3mn\cdot y=-6mny.
\]

\pause

Todėl
\[
9m^2n^2-6mny+y^2
=(3mn)^2-2\cdot3mn\cdot y+y^2
=(3mn-y)^2.
\]

\end{frame}

\begin{frame}{Ne kiekvienas trinaris yra pilnasis kvadratas}

Pavyzdžiui:
\[
x^2+5x+9.
\]

\pause

Pirmasis ir paskutinis nariai yra kvadratai:
\[
x^2=x^2,
\qquad
9=3^2.
\]

\pause

Tačiau
\[
2\cdot x\cdot3=6x,
\]
o vidurinis narys yra $5x$.

\pause

Todėl
\[
x^2+5x+9\neq(x+3)^2.
\]

\end{frame}

\begin{frame}{Pilnojo kvadrato sudarymas}

Kartais trūksta tik paskutinio nario, kad trinaris taptų
pilnojo kvadrato trinariu.

\pause

Tuomet reikiamą narį \textbf{pridedame ir iš karto atimame}.

\pause

\begin{block}{Kodėl taip galima?}
Reiškinio reikšmė nepasikeičia, nes
\[
k-k=0.
\]
\end{block}

\end{frame}

\begin{frame}{Pavyzdys 6}

Išskirkite pilnąjį kvadratą:
\[
x^2+8x.
\]

\pause

Vidurinis narys turi būti
\[
2\cdot x\cdot b=8x,
\]
todėl
\[
b=4.
\]

\pause

Pridedame ir atimame $4^2=16$:
\[
x^2+8x=x^2+8x+16-16.
\]

\pause

\[
x^2+8x=x^2+2\cdot x\cdot4+4^2-16.
\]

\pause

\[
x^2+8x=(x+4)^2-16.
\]

\end{frame}

\begin{frame}{Pavyzdys 7}

Išskirkite pilnąjį kvadratą:
\[
x^2-10x+7.
\]

\pause

Viduriniam nariui $-10x$ gauti reikia $b=5$, nes
\[
-2\cdot x\cdot5=-10x.
\]

\pause

Pridedame ir atimame $5^2=25$:
\[
x^2-10x+7=x^2-10x+25-25+7.
\]

\pause

\[
x^2-10x+7=x^2-2\cdot x\cdot5+5^2-25+7.
\]

\pause

\[
x^2-10x+7=(x-5)^2-18.
\]

\end{frame}

\begin{frame}{Pavyzdys 8}

Išskirkite pilnąjį kvadratą:
\[
4x^2+12x+1.
\]

\pause

Pirmasis narys yra
\[
4x^2=(2x)^2.
\]

\pause

Viduriniam nariui gauti:
\[
2\cdot2x\cdot b=12x,
\]
todėl $b=3$.

\pause

Pridedame ir atimame $3^2=9$:
\[
4x^2+12x+1=4x^2+12x+9-9+1.
\]

\pause

\[
4x^2+12x+1=(2x)^2+2\cdot2x\cdot3+3^2-9+1.
\]

\pause

\[
4x^2+12x+1=(2x+3)^2-8.
\]

\end{frame}

\begin{frame}{Pavyzdys 3.6}

Išskirkite pilnąjį kvadratą:
\[
x^2+3x+1.
\]

\pause

Sprendimas

Vidurinio nario koeficiento pusė yra $\frac{3}{2}$, todėl pridedame ir atimame $\left(\frac{3}{2}\right)^2$:

\[
\begin{aligned}
x^2+3x+1
&=x^2+2\cdot x\cdot\frac{3}{2}+\left(\frac{3}{2}\right)^2-\left(\frac{3}{2}\right)^2+1 \\
&=\left(x+\frac{3}{2}\right)^2-\frac{9}{4}+1 \\
&=\left(x+\frac{3}{2}\right)^2-\frac{5}{4}.
\end{aligned}
\]

\pause

Atsakymas: $\left(x+\frac{3}{2}\right)^2-\frac{5}{4}$.

\end{frame}

\begin{frame}{Pavyzdys 10}
\small

Išskirkite pilnąjį kvadratą:
\[
2x^2+\frac15x+3.
\]

\pause

Iškeliame $2$ prieš pirmuosius du narius:
\[
2x^2+\frac15x+3
=
2\left(x^2+\frac1{10}x\right)+3.
\]

\pause

Skliaustuose vidurinio nario koeficiento pusė yra $\frac1{20}$.

\pause

\[
\begin{aligned}
2x^2+\frac15x+3
&=2\left(
x^2+2\cdot x\cdot\frac1{20}
+\left(\frac1{20}\right)^2
-\left(\frac1{20}\right)^2
\right)+3.
\end{aligned}
\]

\pause

\[
2x^2+\frac15x+3
=
2\left(
\left(x+\frac1{20}\right)^2-\frac1{400}
\right)+3.
\]

\pause

\[
2x^2+\frac15x+3
=
2\left(x+\frac1{20}\right)^2+\frac{599}{200}.
\]

\end{frame}

\begin{frame}{Pavyzdys 11}

Išskirkite pilnąjį kvadratą:
\[
25a^2b^2+10abc+a.
\]

\pause

Pirmasis narys yra
\[
25a^2b^2=(5ab)^2.
\]

\pause

Kad gautume pilnojo kvadrato trinarį, reikia pridėti $c^2$:
\[
2\cdot5ab\cdot c=10abc.
\]

\pause

Todėl
\[
25a^2b^2+10abc+a
=25a^2b^2+10abc+c^2-c^2+a.
\]

\pause

\[
25a^2b^2+10abc+a
=(5ab)^2+2\cdot5ab\cdot c+c^2-c^2+a.
\]

\pause

\[
25a^2b^2+10abc+a
=(5ab+c)^2-c^2+a.
\]

\end{frame}

\begin{frame}{Santrauka}

\begin{block}{Pilnojo kvadrato atpažinimas}
\[
a^2+2ab+b^2=(a+b)^2,
\]
\[
a^2-2ab+b^2=(a-b)^2.
\]
\end{block}

\pause

\begin{block}{Pilnojo kvadrato sudarymas}
Jei trūksta reikiamo nario, jį pridedame ir atimame:
\[
x^2+2bx
=x^2+2bx+b^2-b^2
=(x+b)^2-b^2.
\]
\end{block}

\end{frame}

\end{document}